\documentclass[11pt,a4paper]{article}

\newcommand{\docshorttitle}{HomeSense --- Defence Notes}
% Shared preamble for the HomeSense project documents.
% Compiled with Tectonic (XeTeX). See build.ps1.

\usepackage[a4paper,margin=25mm,headheight=15pt]{geometry}
\usepackage{fontspec}
\usepackage{amsmath}
\usepackage{amssymb}   % provides \square, used for checklist items
\usepackage{microtype}
\usepackage{booktabs}
\usepackage{tabularx}
\usepackage{longtable}
\usepackage{array}
\usepackage{xcolor}
\usepackage{graphicx}
\usepackage{enumitem}
\usepackage{fancyhdr}
\usepackage{titlesec}
\usepackage{caption}
\usepackage{listings}
\usepackage{tikz}
\usetikzlibrary{positioning,arrows.meta,calc,shapes.geometric,fit,backgrounds}
\usepackage[hidelinks]{hyperref}

% --- typography -------------------------------------------------------------
\setmainfont{Times New Roman}[
  BoldFont={Times New Roman Bold},
  ItalicFont={Times New Roman Italic},
  BoldItalicFont={Times New Roman Bold Italic}
]
\setmonofont{Consolas}[Scale=0.86]

\linespread{1.08}
\setlength{\parskip}{0.6em}
\setlength{\parindent}{0pt}
\setlist{itemsep=0.15em,topsep=0.35em,parsep=0pt}

% --- colours ----------------------------------------------------------------
\definecolor{rule}{HTML}{9AA5B1}
\definecolor{codebg}{HTML}{F5F7F9}
\definecolor{codekw}{HTML}{0B5FA5}
\definecolor{codestr}{HTML}{1B7F3B}
\definecolor{codecom}{HTML}{6B7280}
\definecolor{accent}{HTML}{00696D}
\definecolor{boxbg}{HTML}{F0F4F6}

% --- headings ---------------------------------------------------------------
\titleformat{\section}{\large\bfseries}{\thesection}{0.7em}{}
\titleformat{\subsection}{\normalsize\bfseries}{\thesubsection}{0.7em}{}
\titleformat{\subsubsection}{\normalsize\itshape}{\thesubsubsection}{0.7em}{}
\titlespacing*{\section}{0pt}{1.8em}{0.6em}
\titlespacing*{\subsection}{0pt}{1.2em}{0.4em}

% --- headers and footers ----------------------------------------------------
\pagestyle{fancy}
\fancyhf{}
\renewcommand{\headrulewidth}{0.4pt}
\fancyhead[L]{\small\textsc{\docshorttitle}}
\fancyhead[R]{\small SCS 3311}
\fancyfoot[C]{\small\thepage}

% --- tables -----------------------------------------------------------------
\renewcommand{\arraystretch}{1.25}
\newcolumntype{L}[1]{>{\raggedright\arraybackslash}p{#1}}
\captionsetup{font=small,labelfont=bf,skip=6pt}

% --- inline code ------------------------------------------------------------
% Handles underscores without manual escaping, so field names such as
% max_on_duration can be written naturally.
\newcommand{\code}[1]{\texttt{\small #1}}

% --- code listings ----------------------------------------------------------
\lstdefinestyle{project}{
  backgroundcolor=\color{codebg},
  basicstyle=\ttfamily\scriptsize,
  keywordstyle=\color{codekw}\bfseries,
  stringstyle=\color{codestr},
  commentstyle=\color{codecom}\itshape,
  numberstyle=\tiny\color{codecom},
  breaklines=true,
  breakatwhitespace=true,
  captionpos=b,
  frame=single,
  rulecolor=\color{rule},
  framesep=5pt,
  xleftmargin=8pt,
  xrightmargin=4pt,
  showstringspaces=false,
  tabsize=2,
  columns=fullflexible,
  keepspaces=true,
  upquote=true,
}
\lstset{style=project}

\lstdefinelanguage{Kotlin}{
  keywords={package,as,typealias,class,this,super,val,var,fun,for,while,do,null,
    true,false,is,in,throw,return,break,continue,object,if,else,when,try,catch,
    finally,for,do,while,interface,data,sealed,override,private,public,internal,
    suspend,operator,companion,by,import,enum,const},
  sensitive=true,
  comment=[l]{//},
  morecomment=[s]{/*}{*/},
  morestring=[b]",
}

% --- callout box ------------------------------------------------------------
\newcommand{\keypoint}[1]{%
  \par\vspace{0.4em}%
  \noindent\begin{tikzpicture}
    \node[draw=rule,fill=boxbg,line width=0.4pt,rounded corners=2pt,
          inner sep=9pt,text width=\dimexpr\textwidth-24pt\relax,align=justify]
      {#1};
  \end{tikzpicture}%
  \par\vspace{0.4em}%
}

% --- title page -------------------------------------------------------------
\newcommand{\titlepagestandard}[2]{%
\begin{titlepage}
\centering
\vspace*{2.2cm}

{\Large\scshape University of Colombo School of Computing\par}
\vspace{0.4cm}
{\large SCS 3311 --- Mobile Application Design and Development\par}
\vspace{0.3cm}
{\large Mini-Project\par}

\vspace{2.4cm}
{\rule{\textwidth}{0.8pt}}
\vspace{0.7cm}

{\LARGE\bfseries HomeSense\par}
\vspace{0.5cm}
{\Large Smart Home Monitoring and Control System\par}
\vspace{0.7cm}
{\large #1\par}

\vspace{0.7cm}
{\rule{\textwidth}{0.8pt}}

\vspace{2.4cm}

{\large\bfseries Group Members\par}
\vspace{0.6cm}

\begin{tabular}{@{}L{4.6cm}L{3.1cm}L{6.6cm}@{}}
\toprule
\textbf{Name} & \textbf{Index number} & \textbf{Module} \\
\midrule
R.L. Weerasinghe & 23002204 & Synchronisation, simulator and safety worker \\
W.T. Mahagamage  & 23001038 & Device profiles, control interface and reporting \\
D.M. Isakya      & 23000643 & Floor representation and grid mapping \\
\bottomrule
\end{tabular}

\vfill
{#2\par}
\vspace{0.8cm}
\end{titlepage}
}


% Question and answer pairs are the whole structure of this document.
\newcounter{qnum}[section]
\newcommand{\question}[1]{%
  \refstepcounter{qnum}%
  \par\vspace{0.9em}%
  \noindent\textbf{\theqnum.\ #1}\par\vspace{0.25em}%
}
\newcommand{\shared}[1]{%
  \par\vspace{0.7em}\noindent\textbf{#1}\par\vspace{0.25em}%
}

\begin{document}

\titlepagestandard{Defence Notes}{\large August 2026}

\section*{How to use this document}
\addcontentsline{toc}{section}{How to use this document}

Assessment includes individual defence: each member must answer questions on
their own module. This document covers the five questions most likely to be
asked of each member, with an answer for each, followed by questions that span
modules and may be put to anyone.

Read your own section aloud until the answers can be given without reference to
the page.

\section{Questions any member may be asked}

\shared{Why three fields rather than one boolean?}

A single boolean cannot distinguish between a state that has been requested, a
state that is actually in effect, and the system's confidence in that
information. \code{desiredState} records the application's intent,
\code{reportedState} records the hardware's confirmation, and \code{status} is
the server's reconciliation of the two. This is what allows the error and
disconnected states to represent observations rather than assumptions.

\shared{Where does the safety cut-off run, and why there?}

In the worker process, on the server side. Implemented in the application it
would only protect the house while the application was running, which defeats
the purpose of a safety feature.

% ============================================================================
\newpage
\section{R.L. Weerasinghe (23002204) --- synchronisation, simulator and safety}

Owns \code{data/remote/}, \code{data/repository/}, \code{simulator/},
\code{worker/} and \code{database.rules.json}.

\question{Why Realtime Database rather than Firestore?}

Two reasons. Realtime Database provides per-child listeners, so toggling one
slot produces a single callback carrying that device rather than a
retransmission of the whole devices subtree, which matters for a grid of
independently toggled slots. It also provides \code{onDisconnect()}, which lets
the server record a client's disappearance, so the disconnected state is an
observation rather than an inference made independently by each observer.

Firestore's document-level granularity and higher write latency suit a different
access pattern. The cost is the absence of query support and a denormalised
tree, which is acceptable because the tree is small, fixed and documented.

\question{How is the absence of manual refresh actually guaranteed?}

There is no refresh method to call. \code{HomeRepository} exposes no
\code{refresh()}, \code{reload()} or \code{fetch()} operation. The device list is
a \code{StateFlow} supplied by database child listeners, so a change from the
simulator, the worker or another client arrives as a callback and triggers
recomposition.

The repository tests demonstrate this by emitting a change on the source and
asserting the flow carries it, without calling into the repository at all. It is
a property of the design rather than a feature added to it.

\question{Can a client actually not misreport a device's status?}

No, and the mechanism is worth stating precisely. In Realtime Database a write
grant propagates to all descendants of the node where it is declared and cannot
be withdrawn lower down, so a nested \code{.write: false} would have no effect.
Validation rules, however, do restrict client writes, and the Admin SDK bypasses
validation. The rules use that asymmetry: \code{status}, \code{link},
\code{onSince} and \code{mismatchSince} carry validators that accept a write only
when the value is unchanged. A client therefore cannot publish a status value,
clear an error, or suppress a disconnected state, while the worker writes them
freely.

The limitation is that the application and the simulator both authenticate
anonymously, so the rules cannot distinguish them, and a modified client could
write \code{reportedState} as though it were hardware. Separating the roles would
require custom claims and a service to issue them. What no client can do is
write \code{status}, which is what the safety argument depends on.

\question{Why is the cut-off timer not implemented with \code{setTimeout}?}

Because it would be lost when the process terminated, and lost without any
indication. The \code{onSince} value is held in the database and elapsed time is
recomputed from it at every evaluation, so a worker that crashes or is
redeployed resumes correctly: the next pass observes an appliance forty-seven
seconds into a thirty-second limit and switches it off immediately.

There is a test for that scenario and another confirming the cut-off fires at
exactly \code{max\_on\_duration}. Two triggers drive the same pure function ---
child listeners for immediate reaction and a thirty-second sweep for recovery
--- so they cannot produce inconsistent results.

\question{Why does the schedule rule act on boundaries rather than enforcing the
interval?}

Because enforcing the interval would override manual control. A user switching a
scheduled light off at 19:00 would find it on again within a minute. Acting only
at the on-minute and off-minute allows an override to persist until the next
boundary, which matches what a timer is normally expected to do.

Application is idempotent, since a write occurs at a boundary minute only when
the desired state differs from the target, so the sweep cannot apply a
transition twice. The consequence is that a boundary occurring entirely while
the worker is down is not applied retrospectively, which is stated in the report
rather than left implicit.

% ============================================================================
\newpage
\section{W.T. Mahagamage (23001038) --- device profiles and reporting}

Owns \code{ui/device/}, \code{ui/report/}, \code{UsageCsv} and the Room
aggregation queries.

\question{Why is \code{max\_on\_duration} a property of a slot rather than of a
device type?}

Because an iron is connected to an ordinary socket and a scheduled lamp may be
wired to one channel of a gang box. Modelling \code{HAZARD} and \code{LIGHT} as
device kinds would make neither arrangement expressible without a separate type
for every combination. A slot is a controllable point; the appliance is assigned
to it, and both the appliance's properties and its policies belong at that
point. The specification's phrasing, ``specialised slots assigned to
appliances'', indicates the same structure.

\question{The toggles are optimistic. Is that not misleading, given that only
the server knows the actual state?}

It would be if nothing followed. The switch changes position immediately because
waiting for a round trip makes the interface feel unresponsive, but the
application writes only \code{desiredState} and cannot operate a relay. The
optimistic position is held for at most four seconds, and if
\code{reportedState} has not followed, the control reverts and the user is
informed.

The four seconds is chosen to be longer than a normal round trip and shorter
than the worker's ten-second error threshold, so the user is warned shortly
before the status indicator changes.

\question{Why aggregate usage in SQL rather than over the database snapshot?}

Two reasons. It works without a network connection, since Room is local, so the
reporting screen is populated offline. And it scales: summing thousands of
events in application code would require holding the entire log in memory and
recomputing on every emission, whereas SQL performs the aggregation in the
database against an index.

The query sums \code{durationSec} from off and cut-off events only, since those
are the events that close an on-period, so no state reconstruction is required.

\question{Why not derive usage by comparing successive state snapshots?}

Because that would omit every transition occurring while the application was
closed, which is exactly the period the safety worker covers. An appliance
switched on at 14:02 and cut off at 14:32 with the phone unavailable must appear
in the report, and it does, because the worker logged those events when they
occurred.

Usage is append-only and each component logs its own transitions. Comparison
would also double-count after a reconnection; upserting on the database push key
makes replay idempotent.

\question{Why is the energy figure described as an estimate?}

Because it is one. It multiplies the appliance's rated wattage by its on-time,
which assumes full rated draw throughout. A kettle does not behave that way: it
draws heavily while heating and very little afterwards. Measuring properly would
require a metering device reporting actual consumption, which is outside the
scope of a simulation.

Labelling it an estimate wherever it appears is the accurate description, and
accepting a measured value instead would be a one-line change.

% ============================================================================
\newpage
\section{D.M. Isakya (23000643) --- floor representation}

Owns \code{ui/plan/}, \code{ui/floors/}, \code{GridMapper} and
\code{FloorPlanSpec}.

\question{Why store grid cells rather than pixel coordinates?}

Pixel coordinates are not meaningful across devices: the same offset falls in a
different room on a tablet than on a small phone, and rotating the screen would
move every device. A cell is a stable identity --- cell $(4,2)$ is the same
location in the house regardless of what is rendering it.

Translation to pixels happens once, at draw time, from the current canvas size.
It also means the database stores information about the house rather than about
one particular screen.

\question{Why is \code{GridMapper} a separate class rather than arithmetic
inside the composable?}

Because it is the part that can be wrong in a way that is not visible on
inspection. Compose code requires a device or a preview to exercise, whereas a
pure class with no Android dependency runs as a JVM unit test in milliseconds.

There are sixteen tests: the round trip for every cell, taps on the letterbox
margins, taps exactly on a cell boundary and on the plan edge, a zero-sized
canvas, and the property that a given cell resolves identically under phone and
tablet geometry. Inline, none of that would be tested, and a placement error
would first appear during the demonstration.

\question{What is the letterboxing, and why does it matter?}

The plan has a fixed aspect ratio and the available canvas generally does not,
so the plan is fitted within the canvas with its ratio preserved, leaving
margins on two sides. All cell arithmetic is then performed relative to the
fitted rectangle rather than the raw canvas bounds.

Without that, the grid would stretch with the window and devices would drift out
of their rooms as the layout changed. \code{cellAt()} also returns null for a
tap on the margins, which represents ``not on the plan'' rather than ``outside
the grid''.

\question{Why are plans drawn as geometry rather than supplied as images?}

Rendering stays sharp at any screen density with no bitmap scaling; the aspect
ratio is exact, which makes the fitting calculation exact; and because the plans
are original work, there is no third-party image licence to attribute.
\code{FloorPlanSpec} declares rooms, doorways and windows in unit-less
coordinates, and Compose \code{Canvas} renders them.

\question{How does the plan represent a gang box with a fault on one channel?}

The marker shows the most severe of its slot statuses, not an average and not
the first. A fault on one channel must not be concealed by two working channels,
so a three-gang box with one error is drawn as an error marker.

Shape carries the device kind --- circle, square with one indicator per slot,
camera outline --- and fill pattern carries the status: solid, outline, cross,
dashed. Nothing depends on colour alone, which matters both for accessibility
and for legibility when projected.

% ============================================================================
\newpage
\section{Further questions to prepare}

\shared{What happens if two people toggle the same switch simultaneously?}

The last write to \code{desiredState} takes effect, and both clients converge on
whatever the hardware subsequently reports. No locking is required, because the
relay's report rather than either client's intent is the source of truth for
\code{status}.

\shared{Why does the application still show disconnected for a device whose
relay is on?}

Because that is not known. The heartbeat has stopped, so the last reported state
is stale by an unknown interval, and reporting on would present an assumption as
an observation. The \code{onSince} value is retained across a disconnection, so
that if the node returns while still running, the cut-off counts from when it
actually started rather than from the reconnection.

\shared{Room was not required by the specification. Why is it used?}

For offline availability and for usage aggregation, and the report states this
rather than implying it was required. Live device status is not cached, since an
outdated on indicator is more misleading than an honest disconnected one.

\shared{What happens if the worker is not running?}

No status is derived, no cut-off occurs and no schedule is applied. This is a
genuine single point of failure and it is documented as one. The simulator
displays the worker's own heartbeat and reports when it is not running.

\shared{Where does the application write \code{status}?}

It does not, anywhere. \code{RemoteHomeSource}, which defines the complete set of
operations the application can perform against the database, contains no method
that writes it. The database rules enforce the same restriction independently.

\end{document}
